\section{Conclusion}
\label{sec:conclusion}

\slowlab{} evaluates an agent as the operator of a resource-limited experimental campaign.
The design policy determines what evidence can exist, the history reader determines how that
evidence becomes a recommendation, and the evaluator measures both without leaking hidden
truth into the episode. This separation turns a vague failure to ``reason scientifically''
into testable allocation, evidence-use, and recommendation questions.

The current experiments support a narrower conclusion than the Version~1 mechanism story.
Allocation and reading are empirically separable but task dependent. Standard design and
inference aids do not reliably improve endpoint regret in the tested primary conditions; a
fresh high-noise extension preserves that null result, and one lower-noise condition shows
harm. Within-cycle access changes behaviour without improving the final recommendation in
the current synchronous action space. The next decisive tests are matched cross-model
comparisons and action spaces in which interim evidence can change physical capacity.
